\documentclass[11pt,a4paper]{article}

% Focused portfolio report: the implementation is the source of truth.
\usepackage[a4paper,margin=23mm,headheight=24pt]{geometry}
\usepackage{amsmath}
\usepackage{amssymb}
\usepackage{booktabs}
\usepackage{xcolor}
\usepackage{fancyhdr}
\usepackage{hyperref}

% The cached compiler lacks the smaller Latin Modern metrics; keep equation
% subscripts at the document size rather than requesting unavailable 8pt fonts.
\DeclareMathSizes{10.95}{10.95}{10.95}{10.95}

\definecolor{navy}{HTML}{102A43}
\definecolor{bodyblue}{HTML}{243B53}
\definecolor{muted}{HTML}{486581}
\definecolor{ruleblue}{HTML}{D9E2EC}
\definecolor{boxblue}{HTML}{F0F4F8}

\hypersetup{
  colorlinks=true,
  linkcolor=navy,
  urlcolor=navy,
  pdftitle={IAM Request Demand and Capacity Optimization},
  pdfauthor={IAMRequestDemand_capacity_optimization}
}

\pagestyle{fancy}
\fancyhf{}
\lhead{\textcolor{muted}{IAMRequestDemand\_capacity\_optimization}}
\rhead{\textcolor{muted}{Technical report}}
\cfoot{\textcolor{muted}{\thepage}}
\renewcommand{\headrulewidth}{0.4pt}
\renewcommand{\headrule}{\hbox to\headwidth{\color{ruleblue}\leaders\hrule height \headrulewidth\hfill}}

\setlength{\parindent}{0pt}
\setlength{\parskip}{6pt}
\color{bodyblue}

% Keep all headings within the cached 11pt font stack used by the project.
\makeatletter
\renewcommand\section{\@startsection{section}{1}{\z@}%
  {-3.5ex \@plus -1ex \@minus -.2ex}%
  {2.3ex \@plus .2ex}%
  {\normalfont\normalsize\bfseries\color{navy}}}
\renewcommand\subsection{\@startsection{subsection}{2}{\z@}%
  {-3.25ex\@plus -1ex \@minus -.2ex}%
  {1.5ex \@plus .2ex}%
  {\normalfont\normalsize\bfseries\color{bodyblue}}}
\makeatother

\newcommand{\code}[1]{\texttt{\detokenize{#1}}}
\newcommand{\formula}[1]{%
  \begin{center}
    \colorbox{boxblue}{\parbox{0.92\linewidth}{\begin{equation*}#1\end{equation*}}}
  \end{center}}

\begin{document}

\begin{titlepage}
  \centering
  \vspace*{24mm}
  {\normalsize\bfseries\color{navy} IAM Request Demand and Capacity Optimization\par}
  \vspace{6mm}
  {\normalsize A code-aligned probabilistic model for IAM request workload\par}
  \vspace{14mm}
  \fcolorbox{ruleblue}{boxblue}{\begin{tabular}{ccc}
    \textbf{Demand} & \textbf{Workload} & \textbf{Decision}\\
    request count & request-level weights & protected capacity
  \end{tabular}}
  \vfill
  \begin{minipage}{0.84\linewidth}
    \normalsize
    This report documents the current portfolio prototype. It follows the
    implementation path from \code{RequestObservation} to the final capacity
    plan and reports the results of the reproducible synthetic experiments.
    It does not claim production accuracy or vendor-specific behavior.
  \end{minipage}
  \vfill
  {\color{muted}IAMRequestDemand\_capacity\_optimization · version 0.1 · September 2026\par}
\end{titlepage}

\section{Problem and modeling objective}

IAM capacity is often planned from historical averages and operational
experience. That approach can hide two distinct uncertainties: how many
requests will arrive, and how much active backend capacity each request will
consume. The model makes both quantities explicit.

For a future planning bucket $t$, the target is a predictive distribution for
the total request workload $W_t$, not only a point estimate:

\formula{W_t = \sum_{i=1}^{N_t} w_{i,t}}

Here $N_t$ is the future request count and $w_{i,t}$ is the active workload of
request $i$. The capacity policy then protects a selected quantile of $W_t$.
The current prototype uses P95 for the main experiment.

\section{Implementation flow}

The code is deliberately split into small contracts. Data acquisition remains
outside the model, so a future IAM adapter can map vendor records into the same
request observation interface.

\begin{center}
\fbox{\begin{minipage}{0.88\linewidth}\centering
\textbf{IAM adapter} $\rightarrow$ \textbf{RequestObservation}
 $\rightarrow$ \textbf{workload decomposition}\\[2mm]
\textbf{count model}\quad +\quad \textbf{rejection model}\quad +\quad
\textbf{empirical weights}\\[2mm]
\hfill $\rightarrow$ \textbf{compound simulation} $\rightarrow$
 \textbf{capacity plan}
\end{minipage}}
\end{center}

\subsection{Request workload contract}

The request-level workload is implemented in
\code{request_contract.py}. It separates work
that happens before the outcome, approval-workflow execution, and backend tasks
that are conditional on approval:

\formula{w_i = w_{i,\mathrm{pre}} + w_{i,\mathrm{workflow}}
       + \mathbb{1}_{\{\mathrm{approved}\}}w_{i,\mathrm{backend}}}

Backend workload is the sum of resource-time across observed tasks:

\formula{w_{i,\mathrm{backend}} = \sum_j
  \mathrm{active\_seconds}_{ij}\,\mathrm{capacity\_units}_{ij}}

Workflow workload is the active execution used to calculate approval steps,
resolve approvers, and route the request:

\formula{w_{i,\mathrm{workflow}} =
  \mathrm{workflow\_overhead\_seconds}_i\,
  \mathrm{workflow\_capacity\_units}_i}

This sum preserves capacity consumed by parallel tasks. A rejected request
still pays the pre-work and workflow costs; only post-approval backend tasks are
gated by the outcome. Human waiting time is not counted as active capacity
unless the system actually holds a capacity unit while waiting.

\subsection{Historical aggregation and empirical distributions}

\code{workload.py} aggregates observations into period-level request counts and
total workloads. It then fits one weighted empirical distribution for approved
requests and one for rejected requests. No Normal assumption is imposed on
request workload values.

The observations remain visible as actual workload values. Their influence is
controlled by a recency weight:

\formula{a_j = \lambda^j, \qquad 0 < \lambda \leq 1}

Here $j=0$ is the newest observation. With the experiment's $\lambda=0.995$,
old observations are gradually down-weighted without being discarded. The
weighted empirical sampler preserves skew and high-cost tails.

\subsection{Request-count forecast}

\code{forecast.py} fits a Negative Binomial baseline to historical period
counts using method-of-moments estimates. If $\mu$ is the sample mean and $v$
the sample variance, the model uses:

\formula{v = \mu + \phi\mu^2, \qquad
  \widehat{\phi}=\max\left(0,\frac{v-\mu}{\mu^2}\right)}

The $\phi$ term is the dispersion parameter. When observed variance is larger
than the mean, it allows more variation than a Poisson model. Future counts are
sampled through the equivalent Gamma--Poisson mixture: first sample a Poisson
rate from a Gamma distribution, then sample the request count from that rate.

The class also exposes an online correction state:

\formula{s_{t+1}=\omega+\rho s_t+\eta\,\mathrm{score}_t}

The current experiment uses the static fit at each forecast origin; the update
method is the interface for a later continuously updated deployment loop.

\subsection{Rejection model and compound simulation}

The rejection component estimates the probability that a request is rejected
before backend activation. It stores the probability on the log-odds scale,
which keeps the transformed forecast between zero and one. In each simulation,
the flow samples a future request count, samples an outcome for every request,
samples a workload from the matching empirical distribution, and adds all
request workloads into one simulated $W_t$.

\subsection{Capacity decision}

The resulting simulated sample is passed to \code{capacity.py}:

\formula{C_{\mathrm{requests},t}=Q_{0.95}(W_t),\qquad
 C_{\mathrm{spare},t}=\max(0,C_{\mathrm{total},t}-C_{\mathrm{requests},t})}

If protected request capacity exceeds total capacity, the difference is
reported as a deficit. This makes request priority explicit instead of hiding
it inside an average utilization number.

\section{Implementation map}

\begin{center}
\normalsize
\begin{tabular}{p{0.36\linewidth}p{0.51\linewidth}}
\toprule
\textbf{Module} & \textbf{Responsibilities}\\
\midrule
\code{request_contract.py} & Request observations, backend tasks, workload contract\\
\code{workload.py} & Aggregation, recency weights, empirical fit and sampling\\
\code{forecast.py} & Count model, rejection model, compound simulation\\
\code{capacity.py} & Quantile protection, spare capacity, deficit\\
\code{flow.py} & Parameter validation and end-to-end orchestration\\
\code{examples/} & Synthetic data and reproducible experiment runners\\
\bottomrule
\end{tabular}
\end{center}

\section{Experimental design}

The main experiment is synthetic so that the complete flow can be reproduced
without exposing real IAM data. It contains 120 daily periods, approximately
200 requests per period, an approximately 11\% rejection rate, overdispersed
counts, workflow execution for every request, and an 8\% high-cost backend tail
among approved requests.

The capacity experiment uses 10,000 Monte Carlo simulations, a P95 protection
target, total capacity of 3,000 capacity-unit seconds per period, and seed
20260920. The separate walk-forward evaluation uses 84 training periods, 36
held-out periods, and 1,000 simulations at each evaluation origin.

\section{Results}

\subsection{Main capacity experiment}

The generated dataset contains 23,431 request observations. The observed
rejection rate is 10.97\%. The fitted request-count mean is 195.26 requests
per period, and the approved/rejected workload means are 15.47 and 2.70
capacity units per request respectively.

\begin{center}
\begin{tabular}{lr}
\toprule
\textbf{Quantity} & \textbf{Value}\\
\midrule
Simulated total-workload P50 & 2,671.88\\
Simulated total-workload P75 & 3,314.21\\
Simulated total-workload P90 & 3,981.97\\
Simulated total-workload P95 & 4,391.47\\
Simulated total-workload P99 & 5,431.87\\
Configured total capacity & 3,000.00\\
Protected request capacity & 4,391.47\\
Reported capacity deficit & 1,391.47\\
\bottomrule
\end{tabular}
\end{center}

The compact histogram below is a visual summary of the simulated workload
distribution. Bars are scaled relative to the largest bin.

\begin{center}
\setlength{\tabcolsep}{2pt}
\begin{tabular}{c@{\quad}c@{\quad}c@{\quad}c@{\quad}c@{\quad}c@{\quad}c@{\quad}c@{\quad}c@{\quad}c@{\quad}c@{\quad}c}
\rule[-1mm]{6mm}{1mm}&\rule[-1mm]{6mm}{18mm}&\rule[-1mm]{6mm}{51mm}&\rule[-1mm]{6mm}{62mm}&\rule[-1mm]{6mm}{47mm}&\rule[-1mm]{6mm}{27mm}&\rule[-1mm]{6mm}{10mm}&\rule[-1mm]{6mm}{3mm}&\rule[-1mm]{6mm}{2mm}&\rule[-1mm]{6mm}{1mm}&\rule[-1mm]{6mm}{1mm}&\rule[-1mm]{6mm}{1mm}\\
296&944&1592&2240&2888&3536&4184&4832&5480&6128&6776&7424\\
\multicolumn{12}{c}{\normalsize simulated future total-workload bins (capacity units)}
\end{tabular}
\end{center}

\subsection{Walk-forward evaluation}

At each origin, the model sees only the preceding periods, forecasts the next
period, and compares the result with the held-out request count and workload.

\begin{center}
\begin{tabular}{lr}
\toprule
\textbf{Metric} & \textbf{Result}\\
\midrule
Request-count MAE & 60.88 requests\\
Request-count P95 coverage & 91.67\%\\
Workload P50 MAE & 842.99 capacity units\\
Workload P95 coverage & 91.67\%\\
Average-weight MAE & 0.52 units/request\\
Average predicted rejection rate & 11.07\%\\
Average actual rejection rate & 10.59\%\\
\bottomrule
\end{tabular}
\end{center}

The weight estimate is relatively stable in this synthetic scenario, while the
request count is harder to predict. Since total workload is a compound sum,
count uncertainty propagates directly into workload uncertainty.

\subsection{Count-model scenario comparison}

The isolated scenario experiment compares the stationary baseline with a
seasonal extension. The result is diagnostic rather than a model-selection
claim:

\begin{center}
\normalsize
\begin{tabular}{lrrr}
\toprule
\textbf{Scenario} & \textbf{Model} & \textbf{MAE} & \textbf{P95 coverage}\\
\midrule
Stationary & Stationary NB & 61.28 & 94.44\%\\
Stationary & Seasonal NB & 60.02 & 94.44\%\\
Weak seasonal & Stationary NB & 70.84 & 83.33\%\\
Weak seasonal & Seasonal NB & 72.00 & 80.56\%\\
Strong seasonal & Stationary NB & 49.81 & 100.00\%\\
Strong seasonal & Seasonal NB & 38.67 & 100.00\%\\
Random shocks & Stationary NB & 83.74 & 94.44\%\\
Random shocks & Seasonal NB & 82.50 & 91.67\%\\
\bottomrule
\end{tabular}
\end{center}

The comparison illustrates why the current Negative Binomial component should
be treated as a baseline: it can represent overdispersion, but it does not
automatically learn seasonality or external drivers.

\section{Interpretation and limitations}

The portfolio result is not that the synthetic model is already production-
accurate. It is that the capacity question has been decomposed into auditable
components with an explicit risk policy:

\begin{itemize}
  \item request demand is modeled separately from request cost;
  \item request cost is derived from active workflow and backend work;
  \item approval outcomes determine whether backend workload is activated;
  \item simulation converts component uncertainty into a total-workload
        distribution;
  \item the P95 quantile turns that distribution into a request-priority
        capacity decision.
\end{itemize}

The current results are limited by the synthetic data, the pooled count fit,
the method-of-moments baseline, and the absence of a real queue/service-level
model. A production-oriented next step would use aggregated and anonymized IAM
observations, establish a rolling backtest protocol, and add temporal
covariates or an online likelihood-based update only where the data supports
it.

\section{Reproducibility}

From the repository root:

\begin{verbatim}
PYTHONPATH=src python3 examples/run_capacity_experiment.py
PYTHONPATH=src python3 examples/run_walk_forward_evaluation.py
PYTHONPATH=src python3 examples/run_count_model_scenarios.py
\end{verbatim}

The numbers in this report correspond to the seeds and configurations stated
in Section 5. Generated CSV/JSON outputs are intentionally not part of the
portfolio repository; the scripts regenerate them when needed.

\end{document}
